% ─────────────────────────────────────────────────────────────────────────────
% Template «tecnico» — preambolo.
% Incluso dopo \documentclass e dopo metadati.tex, prima di opzioni.tex.
% Non caricare babel e non toccare tocdepth: se ne occupa opzioni.tex, che ha
% l'ultima parola perché rappresenta le scelte fatte nel modulo.
% ─────────────────────────────────────────────────────────────────────────────

% Carattere: pacchetti PSNFSS, non fontspec. La resa non deve dipendere dai
% font installati sulla macchina, e PSNFSS fa parte della distribuzione
% minima: è l'insieme di font su cui si può contare ovunque.
\usepackage[T1]{fontenc}
\usepackage{helvet}
\usepackage{courier}
\renewcommand{\familydefault}{\sfdefault}

\usepackage{geometry}
\geometry{left=22mm,right=22mm,top=20mm,bottom=22mm,headsep=7mm,footskip=12mm}

\usepackage{xcolor}
\definecolor{accento}{HTML}{0090A8}
\definecolor{grafite}{HTML}{16181D}
\definecolor{fumo}{HTML}{5B636E}
\definecolor{nebbia}{HTML}{EDEEEC}

\usepackage{graphicx}
\usepackage{booktabs}
\usepackage{array}
\usepackage{enumitem}
\setlist{itemsep=2pt,parsep=0pt,topsep=4pt}

\usepackage{listings}
\lstset{
  basicstyle=\ttfamily\footnotesize,
  backgroundcolor=\color{nebbia},
  frame=leftline, framerule=1.6pt, rulecolor=\color{accento},
  framesep=6pt, xleftmargin=8pt,
  keywordstyle=\color{accento}\bfseries,
  commentstyle=\color{fumo}\itshape,
  stringstyle=\color{grafite},
  breaklines=true, showstringspaces=false,
  numbers=left, numberstyle=\tiny\color{fumo}, numbersep=8pt
}

\usepackage{titlesec}
\titleformat{\section}{\Large\bfseries\color{grafite}}{\thesection}{0.7em}{}[\vspace{2pt}{\color{accento}\titlerule[1.2pt]}]
\titleformat{\subsection}{\large\bfseries\color{grafite}}{\thesubsection}{0.6em}{}
\titleformat{\subsubsection}{\normalsize\bfseries\color{fumo}}{\thesubsubsection}{0.6em}{}
\titlespacing*{\section}{0pt}{18pt}{8pt}
\titlespacing*{\subsection}{0pt}{12pt}{5pt}

\usepackage{fancyhdr}
\pagestyle{fancy}
\fancyhf{}
\renewcommand{\headrulewidth}{0.4pt}
\renewcommand{\footrulewidth}{0pt}
\fancyhead[L]{\footnotesize\color{fumo}\titolo}
\fancyhead[R]{\footnotesize\color{fumo}\riferimento}
\fancyfoot[L]{\footnotesize\color{fumo}\piede}
\fancyfoot[R]{\footnotesize\color{grafite}\thepage}

\linespread{1.06}
\setlength{\parskip}{5pt}
\setlength{\parindent}{0pt}

\usepackage[colorlinks=true,linkcolor=accento,urlcolor=accento,citecolor=accento]{hyperref}
\hypersetup{pdftitle={\titolo},pdfauthor={\autore}}
